\documentclass[11pt]{report}
% T0_preamble_standalone_A4_En.tex — LOKALER PLATZHALTER
% Im Repo durch die offizielle Korpus-Präambel ersetzen; die
% \input-Zeilen der Wrapper bleiben unverändert gültig.
\usepackage{fontspec}
\usepackage[english]{babel}
\usepackage{csquotes}
\setmainfont{Inter}[BoldFont=Inter-Bold,ItalicFont=Inter-Italic,BoldItalicFont=Inter-BoldItalic]
\setmonofont{JetBrains Mono}[BoldFont=JetBrainsMono-Bold,ItalicFont=JetBrainsMono-Italic]
\usepackage{unicode-math}
\setmathfont{Libertinus Math}
\usepackage{geometry}
\geometry{a4paper,top=2.4cm,bottom=2.4cm,left=2.8cm,right=2.8cm}
\usepackage{microtype,parskip,xcolor}
\usepackage{tcolorbox}
\tcbuselibrary{breakable,skins}
\usepackage{booktabs,longtable,array,amsmath,physics,siunitx}
\sisetup{locale=US}
\usepackage{hyperref}
\hypersetup{colorlinks=true,linkcolor=blue!60!black,urlcolor=green!40!black}

% T0_preamble_patches.tex — LOKALER PLATZHALTER
% Im Repo durch die offizielle Version ersetzen.


\title{Doc.~313 -- No Beginning:\\
The Thermal History on the Time Cycle}
\author{Johann Pascher}
\date{August 2026}

\begin{document}
\maketitle
\tableofcontents

% ============================================================
\chapter*{Doc.~313 -- No Beginning}
\addcontentsline{toc}{chapter}{Doc.~313 -- No Beginning}
% ============================================================

\section*{What this is about}

Doc.~312 ends with obligation~D(ii): can the field configurations
close periodically on the compact time cycle
($\tau_c = 2\pi/H_0 = 91.9$~Gyr) \emph{while} carrying the observed
thermal sequence -- nucleosynthesis, recombination, structure
formation? A circle has no beginning; the observations show a
gradient. This document computes the question
(\texttt{ffgft\_313\_zeitzyklus\_geschichte.py}) and finds three
results and one hard open core:

\begin{enumerate}
  \item \textbf{Antipode finding [K$|$P20]:} the entire observable
        history fills, to $1$--$2\,\%$, exactly \emph{one
        half-cycle}. The last-scattering surface sits at
        $\theta = 0.992\,\pi$, the particle horizon at
        $1.012\,\pi$ -- the \enquote{beginning} is the antipode of
        the time cycle, not an event.
  \item \textbf{Smooth closure [K]:} the observed mass-drift
        profile reaches the antipode with slope
        $\sqrt{\Omega_r} \approx 0.01$ -- the smoothness condition
        of an even periodic closure is met to $1\,\%$, and the
        return half is exactly the photonically unobservable half.
  \item \textbf{Entropy obstruction [K]:} exact $\tau_c$
        periodicity is not dynamically generic (Poincaré
        discrepancy $\sim 10^{10^{104}}$); it must be imposed as a
        boundary condition. Fine-grained admissible,
        coarse-grained open -- the sharpened core of D(ii).
\end{enumerate}

Not part of the A-series; registration in Doc.~190 separately.
References: Doc.~306 ($\partial\text{T}^4 = \emptyset$), 267
(degeneracy), 312 (closure scale, compact time direction, statics
of the relation, torus rest frame).

% ============================================================
\chapter{No boundary, no $t=0$: the starting position}
% ============================================================

FFGFT's position on singularity and big bang stands since
Doc.~306: $\partial\text{T}^4 = \emptyset$ -- the torus has no
boundary, \enquote{what was at $t=0$?} is wrongly posed. Doc.~312
made this concrete and gave it three supports:

\textbf{Geometric:} the time direction is a cycle of length
$91.9$~Gyr. A closed circle has no starting point --
\enquote{$t=0$} is as meaningless on it as \enquote{the beginning
of the equator}.

\textbf{Backwards:} in the mass-drift reading (Doc.~312, statics
of the relation; Wetterich), running back does not send distances
to zero but masses: $\tilde T \cdot m = 1$ then means
$\tilde T \to \infty$ -- no density divergence, but a regime of
ever slower clocks. The \enquote{singularity} is where the clocks
stand still, not where the density diverges.

\textbf{Forwards:} no collapse into a singularity -- the
winding/momentum minimum (appendix of Doc.~312) has $V'' > 0$.

What Doc.~312 left open is the \emph{history}: observations show a
clear thermal sequence. How does a gradient live on a circle? That
is D(ii), and here it is computed, not narrated.

% ============================================================
\chapter{The antipode finding [K$|$P20]}
% ============================================================

\section{Where does the history sit on the cycle?}

Position on the cycle: $\theta = \chi/R_H$, with $\chi$ the light
path distance (the $z \to \chi$ assignment takes the measured
$\Lambda$CDM map; only dimensionless ratios enter, Doc.~267). The
antipode lies at $\theta = \pi$, i.e.\
$\chi = L^{*}/2 = \pi R_H = 14.09$~Gpc. Computed:

\begin{center}
\begin{tabular}{lcc}
\toprule
 & $\chi$ [Gpc] & fraction of $L^{*}/2$ \\
\midrule
Recombination ($z=1090$) & $13.98$ & $0.992$ \\
Particle horizon ($z\to\infty$) & $14.27$ & $1.012$ \\
\bottomrule
\end{tabular}
\end{center}

\begin{center}
\begin{tabular}{lcc}
\toprule
$z$ & $\theta/\pi$ & $m(z)/m_0 = 1/(1{+}z)$ \\
\midrule
$0.5$ & $0.140$ & $0.667$ \\
$1$ & $0.243$ & $0.500$ \\
$3$ & $0.465$ & $0.250$ \\
$10$ & $0.689$ & $0.091$ \\
$100$ & $0.917$ & $0.0099$ \\
$1090$ & $0.992$ & $0.00092$ \\
\bottomrule
\end{tabular}
\end{center}

\textbf{Finding:} the entire observable history -- up to
$z \to \infty$ -- fills one \emph{half-cycle} to $1.2\,\%$. The
last-scattering surface sits at the antipode ($99.2\,\%$), the
formal \enquote{beginning} $1.2\,\%$ behind it. The
\enquote{big bang} is thereby a \emph{place on the cycle} -- the
antipodal region, in the mass-drift reading the region of smallest
masses and slowest clocks -- and not an event before time.

\section{Honest placement of the coincidence [S]}

In $\Lambda$CDM language the same fact is the known
quasi-coincidence $\eta_0 H_0/c \approx \pi$ (here $3.18$ against
$3.14$, i.e.\ $1.2\,\%$) -- there an accident of the measured
$\Omega$ values. In the compact reading it becomes a structural
statement: \emph{the scattering wall sits at the antipode.} For it
to be more than a re-labelled coincidence, the location would have
to follow forward -- e.g.\ from the matter content binding the
wall position to the half-cycle. That is not shown; status [S],
noted as a derivation target.

% ============================================================
\chapter{The smooth closure [K]}
% ============================================================

\section{The profile and its endpoint slope}

Periodicity requires the mass drift $m(\theta)$ to close
continuously on the circle. The simplest closure is an \emph{even}
profile about the antipode: minimum at $\theta=\pi$, symmetric
return on the other half. It is smooth if $dm/d\theta \to 0$ at
the turning point.

In the radiation era $H \simeq H_0\sqrt{\Omega_r}(1{+}z)^2$, hence
analytically

\[
  \eta_0 - \chi = \frac{c}{H_0\sqrt{\Omega_r}}\,\frac{1}{1{+}z}
  \qquad\Longrightarrow\qquad
  \left|\frac{dm}{d\theta}\right|_{\text{end}}
  = \sqrt{\Omega_r} \approx 0.0096.
\]

\textbf{Finding:} the observed profile runs into the endpoint with
slope $\sim 0.01$ -- the smoothness condition of an even periodic
closure is met to $1\,\%$. The deformation required relative to
the measured history is a $1\,\%$ effect at the outermost edge.

\section{Why this is not self-deception}

The return half -- where the mass drift would have to rise again
-- is \emph{exactly} the half behind the scattering wall:
photonically unobservable in principle. That is suspiciously
convenient at first sight, but it is not a choice of this
document; it is a consequence of the antipode finding: the wall
\emph{lies} at $0.992\,\pi$. The closure therefore contradicts no
single photon observation -- and it is nevertheless not untestable
(Ch.~E).

The climate-zone picture from the Doc.~312 discussion thereby
becomes quantitative: the gradient is the map \emph{along} the
cycle (like climate zones along a meridian circle), not an
evolution \emph{out of} a beginning -- and the \enquote{hottest}
point of the map is the antipode.

% ============================================================
\chapter{The entropy obstruction: the core of D(ii) [K]}
% ============================================================

\section{Periodicity is not generic}

The entropy budget of the observable volume is
$S \sim 10^{104}\,k_B$ (dominated by supermassive black holes;
photons: $10^{89}\,k_B$). Poincaré recurrence of a generic state
takes $\sim e^{S/k_B} \sim 10^{10^{104}}$ dynamical times --
against $\tau_c \sim 10^{18.5}$~s. The discrepancy is
$\sim 10^{104}$ \emph{orders of magnitude in the exponent}: exact
$\tau_c$ periodicity cannot be expected dynamically.

\textbf{Consequence:} periodicity is a \textbf{boundary
condition} -- a selection of periodic solutions from the solution
space, not an outcome of evolution. It belongs to the same
category as the choice of the T$^4$ topology itself ([STIPULATION]
in the sense of Doc.~311) and must be booked as such.

\section{Fine- and coarse-grained, separated}

\textbf{Fine-grained [K]:} under unitary evolution the
fine-grained entropy is constant; nothing microscopic forbids an
exactly periodic solution. The boundary condition is admissible.

\textbf{Coarse-grained [open]:} the second law is a statement
about coarse-grained entropy relative to an observer. On the
return half it must decrease -- which is observer-relative
possible (reversal of the correlation direction; no one
\emph{there} would see a violation, since the thermodynamic arrow
turns with it), but unproven. \textbf{Exactly this is the
sharpened core of D(ii):} not the field periodicity (which is
consistent as a boundary condition, Ch.~C), but the proof that a
coarse-grained entropy closure is compatible with the observed
structure formation. Until then: open, without embellishment.

% ============================================================
\chapter{Hawking radiation on the time cycle}
% ============================================================

Does Hawking radiation fit this picture? Yes -- better than
before, because the compact time direction (Doc.~312) has quietly
supplied its foundation. At the same time it sharpens the entropy
core of D(ii). Both computed
(\texttt{ffgft\_313\_hawking\_zyklus.py}).

\section{One principle, three faces [K]}

The three famous \enquote{temperatures out of nothing} -- Unruh,
Gibbons--Hawking, Hawking -- are three separate miracles with
three derivations in the textbooks. Behind them stands \emph{one}
KMS rule: where the time structure is periodic, an observer reads
temperature, $T = \hbar/(k_B\tau)$:

\begin{center}
\begin{tabular}{lll}
\toprule
System & Period $\tau$ & $T = \hbar/(k_B\tau)$ \\
\midrule
Cosmos (compact time) & $2\pi/H_0 = 2.9\times10^{18}$~s &
  $2.63\times10^{-30}$~K \\
Black hole ($M_\odot$) & $8\pi GM/c^3 = 1.24\times10^{-4}$~s &
  $6.17\times10^{-8}$~K \\
Unruh ($a = g$) & $2\pi c/a = 1.9\times10^{8}$~s &
  $3.98\times10^{-20}$~K \\
\bottomrule
\end{tabular}
\end{center}

Doc.~312 \emph{grounded} this rule for the cosmos: the
Gibbons--Hawking temperature follows from the compactness of the
time direction itself, without a horizon. Hawking radiation is
therefore no foreign body and no extra postulate -- it is the
\emph{local special case of the same rule}.

\section{The membrane gets its thermometer [S]}

The membrane picture (narrative Ch.~04: the hole \emph{is} its
fractal membrane, the radiation comes from the oscillating
surface) so far had no answer to \enquote{why exactly this
temperature?}. Now it has one, in the language of the mass map:
mass compresses the time structure ($\tilde T \cdot m = 1$ -- the
closer to the membrane, the slower the clocks), and at the
membrane the local time structure becomes periodic with
$\tau = 8\pi GM/c^3$. Short period, high temperature:

\begin{center}
\begin{tabular}{ll}
\toprule
Object & $T_H$ \\
\midrule
PBH $10^{12}$~kg & $1.2\times10^{11}$~K \\
Earth & $2.1\times10^{-2}$~K \\
Sun & $6.2\times10^{-8}$~K \\
Sgr~A$^{*}$ & $1.5\times10^{-14}$~K \\
SMBH $10^{9}M_\odot$ & $6.2\times10^{-17}$~K \\
\bottomrule
\end{tabular}
\end{center}

The soap bubble of Ch.~04 has a thermometer, and the thermometer
measures the clock map.

\section{The family ladder [S]}

Where does the ladder meet the cosmos? $T_H = T_\text{GH}$
requires $M = c^3/(4GH_0) = 4.66\times10^{52}$~kg -- and its
horizon lies at

\[
  r_s = \frac{2GM}{c^2} = \frac{R_H}{2} \quad\text{(exactly } 0.500\text{)},
\]

at the same time exactly \emph{half} the critical Hubble mass
$c^3/(2GH_0)$. The temperature ladder runs from micro holes to
the whole without a break: the cosmos is the largest family
member.

\section{What Hawking cannot do [K]}

The honest friction with the cycle: evaporation times against
$\tau_c = 91.9$~Gyr:

\begin{center}
\begin{tabular}{lll}
\toprule
Object & $t_\text{evap}$ & $t_\text{evap}/\tau_c$ \\
\midrule
PBH $10^{12}$~kg & $10^{12.4}$~yr & $10^{1.5}$ \\
Sun & $10^{67.3}$~yr & $10^{56.4}$ \\
SMBH $10^{9}M_\odot$ & $10^{94.3}$~yr & $10^{83.4}$ \\
\bottomrule
\end{tabular}
\end{center}

Only primordial holes below the threshold mass

\[
  M^{*} = \left(\frac{\tau_c\,\hbar c^4}{5120\,\pi G^2}\right)^{1/3}
  = 3.3\times10^{11}~\text{kg} \quad(\sim\text{asteroid})
\]

close via Hawking within one cycle. A solar-mass hole is $56$
orders of magnitude too slow, an SMBH $83$. The FFGFT correction
from Ch.~04, $P = P_\text{std}(1 - \xi\ln(M/M_P))$, amounts to
$0.6$--$1.4\,\%$ and changes nothing in this finding.

\textbf{Consequence for D(ii):} stellar and supermassive holes --
the carriers of the $10^{104}$ entropy budget of Ch.~D -- cannot
disappear via Hawking on the cycle. The return half must dismantle
them by another route. Hawking radiation thus fits the picture
thermodynamically perfectly and at the same time \emph{sharpens}
the coarse-grained core of D(ii): it shows how hopelessly slow the
only known return route would be.

\section{Information: two languages, one statement [S]}

Ch.~04 always said: the radiation correlates with the fractal
membrane structure -- nothing is lost
($S_\text{BH}(M_\odot) = 1.05\times10^{77}\,k_B$ of
correlations). Ch.~D of this document says: fine-grained entropy
is constant under unitary evolution -- periodicity is
microscopically admissible. That is the same statement in two
languages: information preservation locally (membrane) and
globally (cycle) are consistent. What remains open is solely the
\emph{coarse-grained} closure -- which was the core all along.

% ============================================================
\chapter{Test surfaces [S]}
% ============================================================

Photons end at the wall ($0.992\,\pi$). Two messengers originate
at or behind the antipode:

\begin{center}
\begin{tabular}{lll}
\toprule
Messenger & Decoupling & Position \\
\midrule
C$\nu$B (neutrino background) & $z \sim 6\times10^{9}$ &
  $\theta \approx 1.01\,\pi$ \\
primordial GW & nominally $z \sim 10^{25}$ &
  $\theta \approx 1.01\,\pi$ \\
\bottomrule
\end{tabular}
\end{center}

A \emph{non-monotonic} signature there -- the reversal of the mass
drift beyond the antipode -- would be the in-principle test
surface of the closure. Practically beyond present measurement,
but nameable; together with the $4.2'$ dipole offset (Doc.~312,
torus rest frame) these are two concrete signatures of the compact
reading.

% ============================================================
\chapter{The antipode condition as a prediction:
$\Omega_m^{*}$ and the $\kappa$ candidate}
% ============================================================

\section{The inversion}

Ch.~B read the location of the scattering wall as a finding and
honestly booked the coincidence $\eta_0 H_0/c \approx \pi$ as
[S]. But the condition can be \emph{inverted}: if the compact
reading requires the particle horizon to close exactly at the
antipode,

\begin{equation}
  \eta_0 = \pi R_H
  \qquad\Longleftrightarrow\qquad
  \int_0^\infty \frac{dz}{E(z)} = \pi,
\end{equation}

then $\Omega_m$ is no longer a free parameter. The condition
solves uniquely (flat, $\Omega_r = 9.3\times10^{-5}$) to:

\begin{equation}
  \boxed{\;\Omega_m^{*} = 0.325\;}
  \qquad\text{against Planck } 0.315 \pm 0.007
  \;\;\Rightarrow\;\; +1.4\,\sigma.
\end{equation}

The cosmic sector thereby obtains, for the first time, a
\emph{prediction for a $\Lambda$CDM parameter} rather than a
reinterpretation -- falsifiable: sharper $\Omega_m$ measurements
(DESI, Euclid) decide whether $0.325$ holds.

\textbf{Declared inheritance:} the calculation uses the
$\Lambda$CDM \emph{form} of $E(z)$ -- including the $(1{+}z)^3$
term with full $\Omega_m$. The prediction is therefore stated
precisely: \emph{the map must be shaped so that
$\int dz/E = \pi$; in $\Lambda$CDM bookkeeping this reads
$\Omega_m = 0.325$.} What carries the matter share of that
bookkeeping is a separate question (Sec.~F.3). Status hence
[K$|$P20\,$\times$\,$E(z)$ form] -- the same inheritance
structure as in Doc.~309, one level deeper.

\section{Consolidation: obligation B and the coincidence become
one condition}

From $\Omega_m^{*}$ follows $\Omega_\Lambda^{*} = 0.675$ and
hence a structural candidate for the so far open order-one factor
of the closure scale (Doc.~312, obligation~B):

\begin{equation}
  \kappa^{*} = 3\,\Omega_\Lambda^{*} = 2.03
  \qquad\text{against measured } \kappa = 2.12
  \;\;(-4.4\,\%).
\end{equation}

The same consolidation logic as with the $10^{123}$ problem
($\to$~P20): two open items -- obligation~B and the [S]
coincidence of Ch.~B -- are reduced to \emph{one} condition. The
cross-check shows this is not circular: $\kappa = 2.12$ would
require $\Omega_m = 0.294$, and then
$\eta_0 H_0/c = 3.268 \ne \pi$. Antipode exactness and today's
Planck $\kappa$ stand in a genuine $4\,\%$ tension -- test
contact, not a fit.

\section{Placements and one new open item}

\textbf{$\Omega_\Lambda$ loses its substance status:} the
\enquote{dark-energy density} is the FLRW translation of the
closure scale; if the antipode condition holds, it is predicted,
not fitted. \textbf{Hubble tension:} the compact reading takes an
unambiguous side with the CMB-anchored value ($H_0 = 66.8$ from
the $\xi^{10}$ chain against Planck $67.4$); the local SH0ES
excess would then have to be a map effect of the calibration [S].

\textbf{Dark matter: two regimes, one open.} Dark \emph{matter}
has deliberately not been mentioned so far -- and the gap belongs
named. \emph{Galactically} (rotation curves, RAR) FFGFT needs no
particle DM: there the transition scale $a_0$ carries the load
(Doc.~308); this document says nothing about that, correctly so.
\emph{At the background level}, however, the $z\to\chi$ map
silently inherits the CDM share: in $\Lambda$CDM bookkeeping,
$\Omega_m^{*} = 0.325$ consists of baryons ($\approx 0.049$)
plus cold dark matter ($\approx 0.276$). What carries this share
in the compact reading -- actual matter, or a property of the
mass-drift profile $m(\theta)$ that books as matter in FLRW
translation -- is unresolved and is booked as a \textbf{second
new open item}. It is independent of the galactic question and
must not be conflated with it.

\textbf{New open item (R61 discipline):} what sets the monopole
temperature $T_0 = 2.725$~K? Dimensionlessly,
$k_B T_0/(m_e c^2) = 4.6\times10^{-10}$ -- \emph{not} an integer
power of $\xi$ ($\log_\xi = 2.41$). By the R61 standard no
numerology is attempted here: the antipode temperature
(equivalently the photon-baryon entropy ratio) is so far unbound
and is booked as an open question of the cosmic sector.

% ============================================================
\chapter{Status overview}
% ============================================================

\begin{center}
\begin{tabular}{lp{6.8cm}l}
\toprule
Statement & Content & Status \\
\midrule
Antipode finding & history $=$ one half-cycle to $1.2\,\%$;
  scattering wall at $0.992\,\pi$ & [K$|$P20] \\
Coincidence placement & $\eta_0 H_0/c \approx \pi$: structural
  statement only after forward derivation & [S] \\
Smooth closure & endpoint slope $\sqrt{\Omega_r} \approx 0.01$;
  even closure met to $1\,\%$ & [K] \\
Return half & $=$ photonically unobservable half; no conflict
  with photon data & [K] \\
Periodicity & boundary condition (Poincaré discrepancy
  $10^{104}$~dex in the exponent), not dynamics & [K] \\
Entropy & fine-grained admissible; coarse-grained closure open --
  sharpened core of D(ii) & open \\
Hawking & local case of the same KMS rule (GH from Doc.~312);
  ladder meets cosmos at $r_s = R_H/2$ & [K]/[S] \\
Evaporation & $M^{*} = 3.3\times10^{11}$~kg; $M_\odot$: $56$~dex
  too slow -- sharpens D(ii) & [K] \\
Test surfaces & C$\nu$B/GW from the antipode; non-monotonicity as
  signature & [S] \\
$\Omega_m^{*}$ & antipode condition inverted: $0.325$
  ($+1.4\,\sigma$); $E(z)$ form inherited & [K$|$P20$\times$$E$] \\
$\kappa^{*}$ & $= 3\Omega_\Lambda^{*} = 2.03$ ($-4.4\,\%$):
  candidate for obligation~B; B $+$ coincidence $=$ one condition
  & [S] \\
DM (background) & $(1{+}z)^3$ share of the map: carrier
  unresolved (galactic: $a_0$, Doc.~308, separate) & open \\
$T_0$ & $k_BT_0/(m_ec^2)$ not a $\xi$ power: antipode temperature
  unbound & open \\
\bottomrule
\end{tabular}
\end{center}

\medskip
\noindent\textbf{Balance:} after this document, \enquote{no big
bang} means precisely: the \enquote{beginning} is the antipode of
the time cycle -- a place, not an event; the observed history is
one half of a map closable smoothly to $1\,\%$; and the remaining
burden lies not with the fields but with coarse-grained entropy.
D(ii) is thereby not settled but reduced to its hard core.

\medskip
\noindent\textbf{Registration:} entry into Doc.~190 will be done
separately.

\end{document}
